%%% Настройки для приложений
\appendix
% Оформление заголовков приложений ближе к ГОСТ:
\renewcommand{\hdngalign}{\centering}   % выравнивание по центру
\renewcommand{\hdngaligni}{}			% выравнивание по центру
\setlength{\otstuplen}{0pt}
\renewcommand*{\chapnumfont}{}
% Добваить слово 'Приложение' независимо от значения параметра chapstyle
\renewcommand*{\printchaptername}{\appendixname}
% Название приложения с новой строки с отступом \midchapskip
\setlength{\midchapskip}{6pt}
\renewcommand*{\afterchapternum}{\par\nobreak\vskip \midchapskip}
% Нумеровка приложений русскими буквами
\renewcommand{\thechapter}{\Asbuk{chapter}}
% Не выносим подразделы приложений в оглавление
\settocdepth{chapter}


%\chapter{Стандартные префиксы ссылок}\label{app:A}
\noskipchapter{Стандартные префиксы ссылок}\label{app:A}

Общепринятым является следующий формат ссылок: \linebreak
\texttt{<prefix>:<label>}.
Например, \verb+\label{fig:knuth}+; \verb+\ref{tab:test1}+;
\verb+\label={lst:external1}+.
В~таблице \cref{tab:tab_pref} приведены стандартные префиксы для различных
типов ссылок.

\begin{table}[htbp]
    \captionsetup{justification=centering}
    \centering{
        \caption{\label{tab:tab_pref}Стандартные префиксы ссылок}
        \begin{tabular}{ll}
            \toprule
            \textbf{Префикс} & \textbf{Описание} \\
            \midrule
            ch:              & Глава             \\
            sec:             & Секция            \\
            subsec:          & Подсекция         \\
            fig:             & Рисунок           \\
            tab:             & Таблица           \\
            eq:              & Уравнение         \\
            lst:             & Листинг программы \\
            itm:             & Элемент списка    \\
            alg:             & Алгоритм          \\
            app:             & Секция приложения \\
            \bottomrule
        \end{tabular}
    }
\end{table}

Для упорядочивания ссылок можно использовать разделительные символы.
Например, \verb+\label{fig:scheemes/my_scheeme}+ или \verb+\label{lst:dts/linked_list}+.
